\scanpage{23}

\diagram{\(\text{(50)}\)\quad
\begin{tikzcd}[ampersand replacement=\&,column sep=small,row sep=small]
1\arrow[r]\&\mathrm{SL}_{\Lambda}(M)'\arrow[r]\arrow[d,hook]\&
\mathrm{GL}_{\Lambda}(M)^{\pm\prime}\arrow[r,"\det"]\arrow[d,hook]\&
\{\pm1\}\arrow[r]\arrow[d,"\unsure{\mathrm{id}|\Vert|\wr}"']\&1\\
1\arrow[r]\&\operatorname{Aut}^{+}(SP_M)\arrow[r]\&
\underbrace{\operatorname{Aut}(SP_M)^{\pm}}_{\substack{\text{automs}\\\text{qui transforment}\\\text{l'orientation }\mu\text{ en}\\\pm\mu}}\arrow[r]\&
\{\pm1\}\arrow[r]\&1
\end{tikzcd}}

% slanted
\margin{où $\mathrm{SL}_\Lambda(M)'\overset{\mathrm{def}}{=}\mathrm{SL}_\Lambda(M)/{\pm1}$\\
$\mathrm{GL}_\Lambda(M)^{\pm\prime}\overset{\mathrm{def}}{=}\mathrm{GL}_\Lambda(M)^\pm/{\pm1}$}

Notons que l'on a

\[
\begin{aligned}
\operatorname{Rep}^{+}(SP_M)&\simeq\overrightarrow{SA}_M\xleftarrow{\sim}\operatorname{Repsp}(M,e_M)/\{\pm\mathrm{id}_M\}\\
&\hspace{5em}\deleted{\unsure{\longleftarrow(x,y)}}\quad
\overset{\shortparallel}{\{\deleted{\unsure{x}}(x,y)\in M^2\mid x\wedge y=e_M\}}\\
\{\dot{x},\dot{y}\}&\hspace{3em}\longleftarrow\!\vert\quad (x,y)\bmod\{\pm\mathrm{id}_M\}
\end{aligned}\tag*{(51)}
\]

d'où le fait que $\mathrm{SL}_\Lambda(M)'$ opère de façon simplement transitive sur $\operatorname{Rep}^{+}(SP_M)$, donc que $\mathrm{GL}_\Lambda(M)^{\pm\prime}$ opère de façon simpl. transitive sur l'ens des repères $\operatorname{Rep}(SP_M)$.

\deleted{\uline{Lemme} Soit $P=(S,A,F,\inserted{R,R'})$ une structure ``$2$-polyèdrale'' connexe \inserted{au sens strict} \deleted{(avec les axiomes \inserted{toute arête incidente à exactement deux sommets et deux faces, pour tout sommet $s$ incident à une face $f$, il y a exactement deux arêtes incidentes à la fois à $s$ et $f$})} -- donc \deleted{la structure} \inserted{le \unsure{contour|couronne}} associé à un sommet,}

% archive sheet number: 22
